(1) 代入 \(x=Ae^{i\omega t},y=Be^{i\omega t}\)，得到矩阵 \(\begin{pmatrix}\omega_x^2-\omega^2&-2i\Omega\omega\\2i\Omega\omega&\omega_y^2-\omega^2\end{pmatrix}(A,B)^T=0\)。行列式为 \((\omega_x^2-\omega^2)(\omega_y^2-\omega^2)-4\Omega^2\omega^2=0\)，即 \(\omega^4-(\omega_x^2+\omega_y^2+4\Omega^2)\omega^2+\omega_x^2\omega_y^2=0\)。注意科里奥利项在行列式中给正的四倍频率系数。

(2) 令 \(S=\omega_x^2+\omega_y^2+4\Omega^2\)，得到 \(\omega_\pm^2=[S\pm\sqrt{S^2-4\omega_x^2\omega_y^2}]/2\)。任一非零 \(\Omega\) 下振幅满足 \(B/A=(\omega_x^2-\omega^2)/(2i\Omega\omega)\)，因比值为纯虚，坐标运动一般呈有旋向的椭圆偏振，而非在固定方向上的两个纯直线模式。上下分支有 \(\omega_+\omega_-=\omega_x\omega_y\)。

(3) 将第一式乘以 \(m\dot x\)、第二式乘以 \(m\dot y\)，相加后科里奥利交叉项 \(-2m\Omega\dot x\dot y+2m\Omega\dot x\dot y\) 正好抵消，得到 \(dE/dt=0\)，其中 \(E=m(\dot x^2+\dot y^2+\omega_x^2x^2+\omega_y^2y^2)/2>0\)。\(S>0\) 且 \(\omega_x^2\omega_y^2>0\)，上式两根严格正，振荡稳定。

(4) 各向同性取 \(\omega_x=\omega_y=\omega_0\)，两根化为 \(\omega_\pm=\sqrt{\omega_0^2+\Omega^2}\pm\Omega\)，对应相反圆偏振。当 \(\Omega\to0\) 且 \(\omega_x\ne\omega_y\)，\(\omega_\pm\to\max(\omega_x,\omega_y),\min(\omega_x,\omega_y)\)；模态极限应回到各坐标轴而不直接套用含 \(1/\Omega\) 的振幅比。非零 \(\Omega\) 使原本共振的两线性模频率分裂，形成避免交叉的典型情形。
